\section{Step 34 Validate the last required counter}\label{chacha-s34}

Compute the block count by rounding the bounded positive length upward, then form the last required counter in 64 bits. Load the initial counter unsigned. Reject a nonzero high word before allocating the frame or writing output.

\subsection{Write and inspect}
Replace the body of region \texttt{stream\_counter\_check}. Retain its \texttt{\# BEGIN} and \texttt{\# END} comments and the supplied labels outside them.

\Needspace{13\baselineskip}
\MARSCodeLabel{MIPS assembly}
\begin{minted}[samepage=true]{mars-idea}
    daddiu $t0, $a3, 63
    dsrl $t0, $t0, 6
    lwu $t1, 48($a0)
    daddu $t0, $t0, $t1
    daddiu $t0, $t0, -1
    dsrl $t0, $t0, 32
    bnez $t0, stream_counter_error
    nop
\end{minted}

Use recipe \texttt{s34-1} with the cumulative program through this edit. Inspect its named checkpoint before continuing to the temporary or normal return.

\subsection{Expected observation}
A 65-byte request at counter one passes; the same count at ffffffff returns two before stores.

\Needspace{9\baselineskip}
\subsection{Check your understanding}
Is one block at counter ffffffff valid?

\begin{itemize}
\item \textbf{A} No; the largest counter is always reserved.
\item \textbf{B} Yes; the last available block may be used.
\item \textbf{C} Only when the nonce is zero.
\end{itemize}
